% Part: first-order-logic 
% Chapter: axiomatic-deduction 
% Section: deduction-theorem-quantifiers

\documentclass[../../../include/open-logic-section]{subfiles}

\begin{document}

\olfileid{fol}{axd}{ddq}

\olsection{Teorema Deduksi mawa Kuantor}

\begin{thm}[Teorema Deduksi]
\ollabel{thm:deduction-thm-q} Yen $\Gamma \cup \{!A\} \Proves !B$,
mula $\Gamma \Proves !A \lif !B$.
\end{thm}

\begin{proof}
Kita lumaku maneh kanthi induksi miturut dawane !!{derivation} kang
formula pungkasane $!B$ saka $\Gamma \cup \{!A\}$.

Bukti kanggo kasus dhasar induksi padha karo bukti ing
\olref[ded]{thm:deduction-thm}.

Kanggo langkah induktif, upamana maneh yen !!{derivation} kang formula
pungkasane $!B$ saka $\Gamma \cup \{!A\}$ nduweni langkah
pungkasan~$!B$ kang dhasare aturan inferensi. Yen aturan inferensi
kang dadi dhasar iku modus ponens, kita lumaku kaya ing bukti
\olref[ded]{thm:deduction-thm}. Yen aturan kasebut \QR, kita
ngerti yen $!B \ident !C \lif \lforall[x][!D(x)]$ lan !!a{formula}
awujud $!C \lif !D(a)$ muncul luwih dhisik ing !!{derivation}, kanthi
$a$ ora muncul ing~$!C$, $!A$, utawa $\Gamma$. Dadi, kita duwe
\begin{align*}
  \Gamma \cup \{!A\} & \Proves !C \lif !D(a),\\
  \intertext{lan hipotesis induksi lumaku, mula kita duwe}
    \Gamma & \Proves !A \lif (!C \lif !D(a)).\\
  \intertext{Miturut}
  & \Proves (!A \lif (!C \lif !D(a))) \lif ((!A \land !C) \lif !D(a))\\
  \intertext{lan modus ponens, kita entuk}
  \Gamma & \Proves (!A \land !C) \lif !D(a).\\
  \intertext{Amarga syarat eigenvariabel isih dicukupi, kita bisa nambahake langkah marang !!{derivation} iki kanthi dhasar \QR, lan entuk}
    \Gamma & \Proves (!A \land !C) \lif \lforall[x][!D(x)].\\
    \intertext{Kita uga duwe}
    & \Proves ((!A \land !C) \lif \lforall[x][!D(x)]) \lif (!A \lif (!C \lif \lforall[x][!D(x)])),\\
    \intertext{mula miturut modus ponens,}
    \Gamma & \Proves !A \lif (!C \lif \lforall[x][!D(x)]),
\end{align*}
yaiku $\Gamma \Proves !A \lif !B$. Cathetan koreksi sumber OLPL-029
lan OLPL-030: siji kurung panutup kang ilang ing formula implikasi
sadurunge dibalekake supaya formula kasebut imbang; kesimpulan pungkasan
uga dibalekake dadi implikasi saka A menyang B, kaya asil ing baris
sadurunge.

Kasus nalika $!B$ nduweni dhasar aturan \QR, nanging awujud
$\lexists[x][!D(x)] \lif !C$, ditinggalake minangka gladhen.
\end{proof}

\begin{prob}
  Rampungna bukti \olref[fol][axd][ddq]{thm:deduction-thm-q}.
\end{prob}

\end{document}
